\documentclass[12pt,a4paper]{article}

% --- Kodierung und Sprache ---
\usepackage[T1]{fontenc}
\usepackage[utf8]{inputenc}
\usepackage[ngerman]{babel}
\usepackage{lmodern}

% --- Mathematik ---
\usepackage{amsmath,amssymb,amsfonts,amsthm,mathtools}
\usepackage{bm}

% --- Layout ---
\usepackage[a4paper,margin=2.7cm]{geometry}
\usepackage{microtype}
\usepackage{setspace}
\onehalfspacing

% --- Tabellen / Links ---
\usepackage{booktabs}
\usepackage{hyperref}
\hypersetup{
	colorlinks=true,
	linkcolor=blue,
	citecolor=blue,
	urlcolor=blue,
	pdftitle={Normschalen der Hurwitz-Ordnung und Einheitenorbits},
	pdfauthor={Thomas Hoffbauer},
	pdfsubject={Explizite Rechnungen zu Hurwitz-Normschalen, Einheitenwirkung und projektiver Struktur}
}

% --- Grafikeinbindung robust ---
\usepackage{graphicx}
\DeclareGraphicsExtensions{.pdf,.png,.jpg}
\graphicspath{{./}{fig/}{figs/}{images/}{img/}}
\newcommand{\includegraphicssafe}[2][]{%
	\IfFileExists{#2}{\includegraphics[#1]{#2}}{%
		\fbox{\parbox[c][3cm][c]{0.8\linewidth}{\centering Grafik \texttt{#2} nicht gefunden.}}%
	}%
}

% --- Theorem-Umgebungen ---
\theoremstyle{plain}
\newtheorem{satz}{Satz}[section]
\newtheorem{lemma}[satz]{Lemma}
\newtheorem{proposition}[satz]{Proposition}
\newtheorem{korollar}[satz]{Korollar}
\newtheorem{konjektur}[satz]{Konjektur}

\theoremstyle{definition}
\newtheorem{definition}[satz]{Definition}
\newtheorem{beispiel}[satz]{Beispiel}
\newtheorem{problem}[satz]{Problem}

\theoremstyle{remark}
\newtheorem{bemerkung}[satz]{Bemerkung}

% --- Makros ---
\newcommand{\OH}{\mathcal O_H}
\newcommand{\QH}{\mathbb H}
\newcommand{\Z}{\mathbb Z}
\newcommand{\Fp}{\mathbb F_p}
\newcommand{\PP}{\mathbb P}
\newcommand{\nrd}{\operatorname{Nrd}}
\newcommand{\Stab}{\operatorname{Stab}}
\newcommand{\Orb}{\operatorname{Orb}}

\title{\textbf{Normschalen der Hurwitz-Ordnung und Einheitenorbits}\\
	\large Explizite Rechnungen, freie Wirkung und erste Hinweise auf projektive lokale Struktur}
\author{Thomas Hoffbauer}
\date{März 2026}

\begin{document}
	
	\maketitle
	
	\begin{abstract}
		Wir untersuchen die Normschalen der Hurwitz-Ordnung
		\[
		X_p := \{x\in \OH : \nrd(x)=p\}
		\]
		für kleine Primzahlen \(p\), sowie die Wirkung der Einheitengruppe \(\OH^\times\) auf diesen Mengen.
		Explizite Berechnungen für \(p=2,3,5,7\) zeigen eine bemerkenswert starre Struktur:
		Die Linkswirkung der \(24\) Hurwitz-Einheiten auf \(X_p\) ist in allen getesteten Fällen frei,
		alle Linksorbits besitzen daher Größe \(24\), und für die ungeraden Primzahlen \(p=3,5,7\)
		ergibt sich numerisch
		\[
		|\OH^\times\backslash X_p|=p+1,
		\qquad
		|X_p|=24(p+1).
		\]
		Darüber hinaus zeigen Reduktionen modulo kleiner ungerader Moduli \(N\), dass globale ganzzahlige
		und halbzahlig-hurwitzsche Repräsentanten lokal zusammenfallen können, während für \(N>p\) in den
		getesteten Beispielen keine Kollisionen auftreten. Die Zahl \(p+1\) legt eine natürliche Beziehung
		zu \(\PP^1(\Fp)\), Hecke-Nachbarschaften und Bruhat--Tits-artigen lokalen Geometrien nahe.
		Wir formulieren die entsprechenden Konjekturen und ordnen die Resultate als ersten strukturellen
		Baustein eines größeren hurwitzschen Spektral- und Operatorprogramms ein.
	\end{abstract}
	
	\section{Einleitung}
	
	Die Hurwitz-Ordnung \(\OH\) in den Hamilton-Quaternionen besitzt eine reiche arithmetische,
	geometrische und gruppentheoretische Struktur. Bereits die endlichen Normschalen
	\[
	X_p=\{x\in \OH:\nrd(x)=p\}
	\]
	für Primzahlen \(p\) tragen nicht nur die offensichtliche Vier-Quadrate-Geometrie, sondern sind
	zugleich mit der Aktion der Einheitengruppe \(\OH^\times\) versehen. Ziel dieses Artikels ist es,
	diese Normschalen für kleine Primzahlen explizit zu berechnen und die durch die Einheiten induzierte
	Orbitstruktur sichtbar zu machen.
	
	Das überraschende Resultat ist, dass sich in den ersten nichttrivialen Fällen eine sehr starre Form zeigt:
	für \(p=3,5,7\) zerfällt \(X_p\) in genau \(p+1\) freie Orbits der Größe \(24\). Damit erscheint die Zahl
	\(p+1\) nicht nur als Zählgröße, sondern als Kandidat einer tieferen lokalen oder projektiven Parametrisierung.
	Dies legt nahe, dass nach Quotientierung durch Einheiten eine kanonische Nachbarschaftsstruktur sichtbar wird,
	die eher an \(\PP^1(\Fp)\), Hecke-Operatoren und Bruhat--Tits-Geometrie erinnert als an eine rein kombinatorische
	Restklasse.
	
	Die Resultate dieses Artikels sind in zweierlei Hinsicht zu lesen: zum einen als explizite numerische
	Bestandsaufnahme der kleinen Primfälle, zum anderen als konzeptioneller Hinweis auf eine allgemeinere
	Theorie der Hurwitz-Schalen und ihrer lokalen Reduktionen.
	
	\section{Hurwitz-Ordnung und Normschalen}
	
	\begin{definition}[Hurwitz-Ordnung]
		Die Hurwitz-Ordnung \(\OH\subset \QH\) besteht aus allen Quaternionen
		\[
		x=a+bi+cj+dk
		\]
		mit entweder \(a,b,c,d\in \Z\) oder \(a,b,c,d\in \Z+\tfrac12\).
	\end{definition}
	
	\begin{definition}[Reduzierte Norm]
		Für \(x=a+bi+cj+dk\in \QH\) definieren wir die reduzierte Norm durch
		\[
		\nrd(x):=a^2+b^2+c^2+d^2.
		\]
	\end{definition}
	
	\begin{definition}[Normschale]
		Für eine Primzahl \(p\) definieren wir die zugehörige Hurwitz-Normschale als
		\[
		X_p := \{x\in \OH : \nrd(x)=p\}.
		\]
	\end{definition}
	
	\begin{bemerkung}
		Die Menge \(X_p\) enthält sowohl ganzzahlige als auch halbzahlig-hurwitzsche Repräsentanten.
		Explizit ergibt sich für halbzahlig-hurwitzsche Elemente die Bedingung, dass
		\[
		x=\frac{m_1}{2}+\frac{m_2}{2}i+\frac{m_3}{2}j+\frac{m_4}{2}k
		\]
		mit ungeraden ganzen Zahlen \(m_r\) und
		\[
		m_1^2+m_2^2+m_3^2+m_4^2=4p.
		\]
	\end{bemerkung}
	
	\section{Die Einheitengruppe und ihre Wirkung}
	
	\begin{definition}[Hurwitz-Einheiten]
		Die Einheitengruppe der Hurwitz-Ordnung ist
		\[
		\OH^\times := \{u\in \OH : \nrd(u)=1\}.
		\]
	\end{definition}
	
	\begin{lemma}
		Es gilt
		\[
		|\OH^\times|=24.
		\]
	\end{lemma}
	
	\begin{proof}
		Die \(24\) Einheiten bestehen aus
		\[
		\{\pm1,\pm i,\pm j,\pm k\}
		\]
		sowie den \(16\) Quaternionen der Form
		\[
		\frac{\pm1\pm i\pm j\pm k}{2}.
		\]
		Alle diese Elemente haben Norm \(1\), und weitere Norm-\(1\)-Elemente existieren in \(\OH\) nicht.
	\end{proof}
	
	\begin{definition}[Linkswirkung]
		Die Gruppe \(\OH^\times\) wirkt durch Linksmultiplikation auf \(X_p\):
		\[
		\OH^\times\times X_p\to X_p,\qquad (u,x)\mapsto ux.
		\]
	\end{definition}
	
	\begin{lemma}[Invarianz der Normschale]
		Für jede Primzahl \(p\) ist \(X_p\) unter der Linkswirkung von \(\OH^\times\) invariant.
	\end{lemma}
	
	\begin{proof}
		Für \(u\in \OH^\times\) und \(x\in X_p\) gilt
		\[
		\nrd(ux)=\nrd(u)\nrd(x)=1\cdot p=p.
		\]
		Also liegt \(ux\) wieder in \(X_p\).
	\end{proof}
	
	\begin{definition}[Orbit und Stabilisator]
		Für \(x\in X_p\) definieren wir den Linksorbit
		\[
		\OH^\times\cdot x:=\{ux:u\in \OH^\times\}
		\]
		und den Stabilisator
		\[
		\Stab_{\OH^\times}(x):=\{u\in \OH^\times:ux=x\}.
		\]
	\end{definition}
	
	\section{Explizite Berechnung der kleinen Primfälle}
	
	Die vollständige explizite Enumeration der Normschalen \(X_p\) liefert für die ersten Primzahlen:
	
	\[
	|X_2|=24,\qquad |X_3|=96,\qquad |X_5|=144,\qquad |X_7|=192.
	\]
	
	\begin{beispiel}[Die Normschale \(X_2\)]
		Die Norm-\(2\)-Elemente sind genau die Quaternionen vom Typ
		\[
		(\pm1,\pm1,0,0)
		\]
		mit allen Permutationen der Koordinaten. Entsprechend besteht \(X_2\) aus den \(24\) Elementen
		\[
		\pm1\pm i,\quad \pm1\pm j,\quad \pm1\pm k,\quad
		\pm i\pm j,\quad \pm i\pm k,\quad \pm j\pm k.
		\]
	\end{beispiel}
	
	\begin{beispiel}[Ganzzahlige und halbzahlig-hurwitzsche Elemente in \(X_3\)]
		Für \(p=3\) treten sowohl ganzzahlige Repräsentanten vom Typ
		\[
		(\pm1,\pm1,\pm1,0)
		\]
		als auch halbzahlig-hurwitzsche Repräsentanten vom Typ
		\[
		\left(\pm\frac32,\pm\frac12,\pm\frac12,\pm\frac12\right)
		\]
		auf. Insgesamt ergibt dies
		\[
		|X_3|=96.
		\]
	\end{beispiel}
	
	\begin{bemerkung}
		Die Fälle \(p=5\) und \(p=7\) zeigen bereits deutlich, dass die Normschalen nicht auf eine einzige
		einfache Form reduziert werden können. Vielmehr mischen sich ganzzahlige und halbzahlig-hurwitzsche
		Typen, die erst unter der Einheitenwirkung als Teile größerer Orbits verständlich werden.
	\end{bemerkung}
	
	\section{Freie Wirkung und Orbitzerlegung}
	
	Die folgenden Resultate beruhen auf vollständigen expliziten Rechnungen für \(p=2,3,5,7\).
	
	\begin{proposition}[Freie Wirkung in den kleinen Primfällen]
		Für \(p=2,3,5,7\) ist die Linkswirkung von \(\OH^\times\) auf \(X_p\) frei.
	\end{proposition}
	
	\begin{proof}
		Für jedes \(x\in X_p\) wurde der Stabilisator
		\[
		\Stab_{\OH^\times}(x)=\{u\in \OH^\times:ux=x\}
		\]
		vollständig berechnet. In allen Fällen ergab sich
		\[
		\Stab_{\OH^\times}(x)=\{1\}.
		\]
	\end{proof}
	
	\begin{korollar}
		Für \(p=2,3,5,7\) besitzt jeder Linksorbit in \(X_p\) die Größe \(24\).
	\end{korollar}
	
	\begin{proof}
		Dies folgt unmittelbar aus dem Orbit-Stabilisator-Prinzip:
		\[
		|\OH^\times\cdot x|
		=
		\frac{|\OH^\times|}{|\Stab_{\OH^\times}(x)|}
		=
		24.
		\]
	\end{proof}
	
	\begin{proposition}[Orbitzahlen in den kleinen Primfällen]
		Für \(p=2,3,5,7\) besitzt \(X_p\) genau
		\[
		1,\ 4,\ 6,\ 8
		\]
		Linksorbits unter \(\OH^\times\).
	\end{proposition}
	
	\begin{proof}
		Da jeder Orbit nach dem vorigen Korollar Größe \(24\) besitzt, folgt
		\[
		|X_2|/24=1,\qquad |X_3|/24=4,\qquad |X_5|/24=6,\qquad |X_7|/24=8.
		\]
	\end{proof}
	
	\begin{satz}[Numerisch verifizierte Orbitformel]
		Für die getesteten ungeraden Primzahlen \(p=3,5,7\) gilt numerisch
		\[
		|\OH^\times\backslash X_p|=p+1,
		\]
		und damit
		\[
		|X_p|=24(p+1).
		\]
		Für \(p=2\) besteht \(X_2\) aus genau einem Orbit der Größe \(24\).
	\end{satz}
	
	\begin{proof}
		Die Aussage ist die unmittelbare Kombination der vorangehenden Propositionen.
	\end{proof}
	
	\begin{bemerkung}
		Die Orbitrepräsentanten zeigen bereits eine systematische Form. Für \(p=5\) lassen sich die
		sechs Orbits etwa durch Repräsentanten
		\[
		-2\pm i,\qquad -2\pm j,\qquad -2\pm k
		\]
		beschreiben. Für \(p=7\) erscheinen acht Repräsentanten, die sich als Vorzeichenmuster
		eines halbzahlig-hurwitzschen Typs mit Hauptkoordinate \(-\frac52\) lesen lassen.
		Dies deutet darauf hin, dass die \(p+1\) Orbits nicht zufällig auftreten, sondern eine
		verborgene projektive oder lokale Parametrisierung besitzen.
	\end{bemerkung}
	
	\section{Reduktion modulo \texorpdfstring{$N$}{N} und lokale Schalenphänomene}
	
	Um lokale Phänomene sichtbar zu machen, betrachten wir die koordinatenweise Reduktion modulo
	ungerader Zahlen \(N\). Für halbzahlig-hurwitzsche Einträge ist dies sinnvoll, da \(2\) modulo
	ungeradem \(N\) invertierbar ist.
	
	Die expliziten Rechnungen zeigen:
	
	\begin{itemize}
		\item Für \((p,N)=(2,3)\) gilt
		\[
		|X_2 \bmod 3|=24=|X_2|.
		\]
		\item Für \((p,N)=(3,5)\) gilt
		\[
		|X_3 \bmod 5|=96=|X_3|.
		\]
		\item Für \((p,N)=(5,7)\) gilt
		\[
		|X_5 \bmod 7|=144=|X_5|.
		\]
	\end{itemize}
	
	In diesen Beispielen mit \(N>p\) treten also keine Kollisionen auf.
	
	Demgegenüber zeigen kleine Moduli einen starken Kollaps:
	
	\begin{itemize}
		\item Für \((p,N)=(5,3)\) gilt
		\[
		|X_5 \bmod 3|=24.
		\]
		\item Für \((p,N)=(7,5)\) gilt
		\[
		|X_7 \bmod 5|=96.
		\]
	\end{itemize}
	
	\begin{bemerkung}
		Die Beispiele zeigen insbesondere, dass ganzzahlige Repräsentanten wie
		\[
		-2j+k,\qquad j-2k
		\]
		und halbzahlig-hurwitzsche Repräsentanten wie
		\[
		-\frac32-\frac32i-\frac12j-\frac12k
		\]
		modulo kleinem \(N\) auf dieselbe Restklasse fallen können. Dies ist ein klarer Hinweis darauf,
		dass die rohe Bildmenge \(X_p \bmod N\) noch zu fein ist und erst nach Quotientierung durch Einheiten
		oder Doppelklassen ihre eigentliche lokale Struktur preisgibt.
	\end{bemerkung}
	
	\begin{proposition}[Numerische Negationssymmetrie]
		In allen getesteten Fällen ist die reduzierte Bildmenge \(X_p \bmod N\) unter der Abbildung
		\[
		q \longmapsto -q
		\]
		stabil.
	\end{proposition}
	
	\begin{proof}
		Dies wurde in allen explizit berechneten Beispielen numerisch verifiziert.
	\end{proof}
	
	\section{Projektive Interpretation}
	
	Die Zahl \(p+1\) ist in der vorangehenden Orbitformel nicht nur numerisch auffällig, sondern
	arithmetisch hochgradig suggestiv. Es gilt
	\[
	|\PP^1(\Fp)|=p+1.
	\]
	Genau dieselbe Zahl erscheint als natürliche Nachbarschaftszahl in Hecke- und Bruhat--Tits-artigen
	Situationen.
	
	Dies motiviert die folgenden Konjekturen.
	
	\begin{konjektur}[Projektive Parametrisierung]
		Für jede ungerade Primzahl \(p\) ist die Orbitmenge
		\[
		\OH^\times\backslash X_p
		\]
		kanonisch mit der projektiven Geraden
		\[
		\PP^1(\Fp)
		\]
		identifizierbar.
	\end{konjektur}
	
	\begin{konjektur}[Allgemeine Orbitformel]
		Für jede ungerade Primzahl \(p\) wirkt \(\OH^\times\) frei auf \(X_p\), und es gilt
		\[
		|\OH^\times\backslash X_p|=p+1.
		\]
		Insbesondere folgt
		\[
		|X_p|=24(p+1).
		\]
	\end{konjektur}
	
	\begin{konjektur}[Lokale Nachbarschaftsdeutung]
		Die durch Quotientierung von \(X_p\) entstehende Struktur ist nicht bloß kombinatorisch,
		sondern trägt eine kanonische lokale Nachbarschaftsgeometrie, die im geeigneten Sinn mit
		einer Hecke- oder Bruhat--Tits-Struktur korrespondiert.
	\end{konjektur}
	
	\section{Konzeptionelle Einordnung}
	
	Die hier gewonnenen Resultate sind klein, aber strukturell bemerkenswert. Drei Punkte stechen hervor:
	
	\begin{enumerate}
		\item Die Normschalen \(X_p\) sind keineswegs formlos, sondern zerfallen in den getesteten
		Fällen in freie Einheitenorbits fester Größe \(24\).
		
		\item Für ungerade Primzahlen \(p=3,5,7\) ergibt sich numerisch genau die Orbitzahl \(p+1\),
		was unmittelbar auf projektive und lokale arithmetische Strukturen hinweist.
		
		\item Die Reduktionen modulo kleinem \(N\) zeigen, dass global verschiedene ganzzahlige und
		halbzahlig-hurwitzsche Repräsentanten lokal zusammenfallen können. Dies deutet darauf hin,
		dass die kanonische lokale Schale nicht \(X_p\) selbst, sondern ein geeigneter Quotient
		nach Einheiten oder Doppelklassen ist.
	\end{enumerate}
	
	Damit wird ein möglicher Weg sichtbar:
	\[
	X_p
	\;\longrightarrow\;
	\OH^\times\backslash X_p
	\;\longrightarrow\;
	\text{lokale Reduktion}
	\;\longrightarrow\;
	\text{projektive bzw. Hecke-artige Nachbarschaft}.
	\]
	
	\section{Ausblick}
	
	Die naheliegenden nächsten Schritte sind:
	
	\begin{enumerate}
		\item die Berechnung der Rechtsorbits und der Doppelorbits
		\[
		\OH^\times\backslash X_p/\OH^\times,
		\]
		\item die Suche nach einer expliziten Parametrisierung der \(p+1\) Orbits durch
		\(\PP^1(\Fp)\),
		\item die Untersuchung, ob und wie die modulare Reduktion der Orbitrepräsentanten
		lokale Nachbarschaftsrelationen reproduziert,
		\item die Einbettung dieser Struktur in ein Hecke- oder Bruhat--Tits-inspiriertes
		Operatorbild.
	\end{enumerate}
	
	Gerade der Zusammenhang zwischen der globalen Hurwitz-Geometrie und der lokalen Zahl \(p+1\)
	scheint der eigentliche Schlüssel zu sein. Die hier dargestellten kleinen Primfälle bilden dafür
	den ersten belastbaren numerischen Ausgangspunkt.
	
	\section{Schluss}
	
	Die expliziten Rechnungen zu den Normschalen \(X_p\) der Hurwitz-Ordnung zeigen bereits in den
	ersten Primfällen eine überraschend starre Struktur. Die freie Wirkung der \(24\) Hurwitz-Einheiten
	auf \(X_p\), die Orbitgröße \(24\) und die numerische Orbitzahl \(p+1\) für \(p=3,5,7\) sprechen dafür,
	dass hinter den Normschalen nicht nur eine Vier-Quadrate-Geometrie, sondern eine projektiv-lokale
	arithmetische Struktur steht.
	
	Der Status dieser Erkenntnisse ist bewusst klar gehalten: Die kleinen Fälle sind vollständig gerechnet,
	die allgemeine Formel bleibt Konjektur. Gerade diese methodische Trennung macht die beobachtete Struktur
	wissenschaftlich interessant. Denn sie zeigt, dass aus expliziten Hurwitz-Rechnungen heraus bereits
	eine Geometrie sichtbar wird, die mit klassischen lokalen und automorphen Mustern in Resonanz steht.
	
\end{document}